\chapter{结论与进一步工作}

\section{主要结论}
本文针对自动化车床刀具问题，建立了“寿命分布 + 检查策略 + 换刀策略 + 成本函数”的统一优化模型，并形成双层求解框架。主要结论如下：
\begin{enumerate}
  \item 问题可等价为单位合格品平均成本最小化的约束非线性优化问题；
  \item 理想判定与误判判定两场景可在同一模型下统一处理，仅需调整质量参数与误停成本项；
  \item 最优检查策略通常呈现“寿命前段较疏、后段较密”的结构特征；
  \item 误判成本较高时，适度复检机制可有效抑制误停损失。
\end{enumerate}

\section{模型价值}
本文模型不仅能输出最优检查节点和换刀策略，还可通过成本分解和参数敏感性分析解释“为什么该策略最优”。这对工业场景中的策略沟通和管理决策具有实际意义。

\section{局限性}
当前模型仍有以下简化：
\begin{itemize}
  \item 刀具寿命分布设为静态，不包含工况漂移；
  \item 误判概率采用固定参数，未考虑检测设备状态变化；
  \item 未引入多刀具并行工位与产线联动约束。
\end{itemize}

\section{后续工作}
后续可沿以下方向拓展：
\begin{enumerate}
  \item 引入在线学习机制，动态更新寿命分布参数；
  \item 基于贝叶斯决策融合多源检测信号，降低误判；
  \item 将模型扩展到多工位协同优化，形成产线级维护调度策略。
\end{enumerate}
